\chapter{Altholz: Regulierung und Klassifizierung}
\label{ch:background}

Bevor die materialphysikalische Frage nach der \gls{restfestigkeit} eines
gebrauchten Balkens behandelt werden kann, ist zu klaeren, unter welchen
rechtlichen und normativen Bedingungen ein solcher Balken ueberhaupt wieder
tragend eingebaut werden darf. Dieses Kapitel ordnet die einschlaegigen
Regelwerke in zwei Straenge: das Abfallrecht, das darueber entscheidet, ob ein
Stueck \gls{altholz} als Baustoff oder als Abfall behandelt wird, und das
Normenwerk des Holzbaus, das darueber entscheidet, mit welchen
Festigkeitswerten ein zugelassenes Bauteil bemessen wird. Keiner der beiden
Straenge, so das Ergebnis des Kapitels, kennt ein Verfahren fuer die
\gls{lastgeschichte} eines Bauteils.

\section{Altholzverordnung und die Kategorien A I bis A IV}
\label{sec:altholzv}

Die \acrfull{altholzv} regelt in Deutschland die Anforderungen an die
Verwertung und Beseitigung von Altholz \parencite{altholzv}. Sie ordnet
Altholz nach dem Grad seiner stofflichen Belastung in vier Kategorien und
knuepft daran die zulaessigen Verwertungswege. Die Recherche fuer diese
Arbeit konnte die Verordnung als Rechtsquelle und ihre Gliederung
nachvollziehen, nicht jedoch den vollstaendigen Text der Anlagen, in denen die
Kategorien im Detail definiert und die analytischen Schwellenwerte festgelegt
sind. Die folgende Darstellung gibt daher den Belastungsgradienten nach Angaben
der Fachliteratur wieder und nicht den Wortlaut der Verordnung:

\begin{description}
  \item[A I] Naturbelassenes Altholz ohne Beschichtungen, Anstriche oder
    Klebstoffanteile.
  \item[A II] Gestrichenes, beschichtetes oder verleimtes Altholz, dessen
    Anstriche und Klebstoffe nach Angaben der Recherche keine giftigen
    Substanzen enthalten.
  \item[A III] Altholz mit Beschichtungen, die Schadstoffe enthalten koennen,
    etwa schwermetallhaltige Anstriche.
  \item[A IV] Mit Holzschutzmitteln behandeltes Altholz, beispielsweise
    kreosot- oder pentachlorphenolhaltiges Material.
\end{description}

Fuer die tragende Wiederverwendung kommt nach Angaben der Recherche allein die
Kategorie A I ohne zusaetzliche Einschraenkungen in Frage. Material der
Kategorie A II ist nur nach einer vorgeschalteten Schadstoffanalytik und unter
engen Voraussetzungen denkbar, weil Beschichtungen sowohl ein stoffliches
Risiko darstellen als auch die visuelle Beurteilung des Holzes verdecken. Die
Kategorien A III und A IV sind fuer eine tragende Wiederverwendung
ausgeschlossen und stehen praktisch nur der stofflichen oder energetischen
Verwertung offen. Welche konkreten Konzentrationsgrenzen dabei gelten,
konnte die Recherche nicht am Normtext verifizieren; die Verordnung legt die
entsprechenden Pruefanforderungen in ihren Anlagen fest, deren Volltext nicht
frei zugaenglich war. Festzuhalten ist, dass die Kategorisierung ausschliesslich stofflich
argumentiert: sie sagt etwas ueber Schadstoffe aus und nichts ueber
Tragfaehigkeit.

\section{Wiederverwendung gegenueber stofflicher Verwertung}
\label{sec:hierarchie}

Der abfallrechtliche Rahmen oberhalb der \acrshort{altholzv} bildet das
\acrfull{krwg}, das die europaeische Abfallrahmenrichtlinie in nationales
Recht umsetzt \parencite{krwg,eu2008waste}. Beide Regelwerke ordnen die
Behandlungsoptionen in einer Hierarchie: Vermeidung steht an erster Stelle,
darauf folgen die Vorbereitung zur Wiederverwendung und die Wiederverwendung
selbst, danach das Recycling, danach die sonstige, insbesondere energetische
Verwertung und zuletzt die Beseitigung. Die exakte Staffelung und Nummerierung der Stufen konnte die Recherche nicht am Gesetzestext pruefen; die Rangfolge
Wiederverwendung vor Recycling vor energetischer Verwertung ist jedoch
durchgaengig belegt.

Die Unterscheidung ist fuer diese Arbeit nicht bloss terminologisch. Bei der
stofflichen Verwertung wird das Holz zerspant und zu Span-, Faser- oder
Sperrholzplatten weiterverarbeitet. Der Werkstoff bleibt erhalten, das Bauteil
und mit ihm die gewachsene Faserstruktur gehen verloren. Genau deshalb ist die
stoffliche Verwertung gegenueber Kontaminationen der Kategorien A II und A III
toleranter: das Material wird ohnehin aufgeloest. Bei der Wiederverwendung
hingegen bleibt der Balken als Bauteil bestehen und damit auch alles, was
Jahrzehnte unter Last mit ihm geschehen ist. Die Abfallhierarchie verlangt,
diesen Weg zu bevorzugen, und stellt damit eine Frage, die sie selbst nicht
beantwortet: mit welcher Tragfaehigkeit darf man rechnen, wenn ein Bauteil
seine Geometrie behaelt, seine Belastungsvorgeschichte aber unbekannt ist. Die
Regulierung erzeugt den Bedarf an genau jener Kenntnis, die diese Arbeit
untersucht.

\section{Sortierung und Festigkeitsklassen}
\label{sec:sortierung}

Die \gls{sortierung} von Bauholz nach Tragfaehigkeit erfolgt entweder visuell
oder maschinell. Die visuelle Sortierung von Nadelschnittholz ist in
DIN 4074 geregelt \parencite{din4074}; die Norm teilt Nadelholz in
Sortierklassen ein, die mit S7, S10 und S13 bezeichnet werden, und stuetzt
sich dabei auf beobachtbare Merkmale: Groesse und Lage der Aeste, die
Faserneigung gegenueber der Stabachse, Risse und Spalten, Faeulnis und
Verfaerbungen sowie die Jahrringbreite als Indikator der
Wuchsgeschwindigkeit. Die Norm setzt dabei sauberes, unbeschichtetes Material
voraus, weshalb sie auf Altholz nur nach einer Vorreinigung und nur mit
Vorbehalt anwendbar ist: Anstriche verdecken Merkmale oder erzeugen
Scheinmerkmale, und verbliebene Naegel oder Schrauben machen eine
Beurteilung unbrauchbar.

Die maschinelle Sortierung regelt EN 14081 \parencite{en14081}. Sie
unterscheidet maschinenkontrollierte von ausgangskontrollierten Systemen,
verlangt eine Typpruefung und eine laufende Werksproduktionskontrolle jeder
Sortiermaschine und schreibt eine begleitende visuelle Nachkontrolle vor, weil
die Maschine bestimmte Merkmale grundsaetzlich nicht erfasst. Ihr Vorteil
gegenueber der visuellen Sortierung liegt in der objektiven Messung der
Steifigkeit, ihre Grenze in verborgenen Fehlstellen wie inneren Rissen oder
Faeulnisnestern.

Beide Wege muenden in die Festigkeitsklassen nach EN 338
\parencite{en338}. Nadelholz wird dort mit C-Klassen von C14 bis C40
bezeichnet, Laubholz mit D-Klassen; die Kennzahl gibt die charakteristische
Biegefestigkeit in Newton je Quadratmillimeter an. Die Klasse C24
beispielsweise ist durch eine charakteristische Biegefestigkeit von
24 N/mm$^2$, eine Zugfestigkeit parallel zur Faser von 14,5 N/mm$^2$, eine
Druckfestigkeit parallel zur Faser von 21 N/mm$^2$ und eine Schubfestigkeit
von 4,0 N/mm$^2$ gekennzeichnet, C30 durch 30, 19, 24 und 4,0 N/mm$^2$. Die
Zuordnung der Sortierklassen zu diesen Festigkeitsklassen leistet EN 1912
\parencite{en1912}, die fuer die gaengigen europaeischen Nadelholzarten S7 auf
C16, S10 auf C24 und S13 auf C30 abbildet.

Die charakteristischen Werte selbst sind nicht gesetzt, sondern abgeleitet.
EN 384 beschreibt das Verfahren, mit dem sie aus Pruefergebnissen an einer
Stichprobe fuer eine definierte Holzpopulation bestimmt werden
\parencite{en384}, und verweist fuer die Pruefungen auf EN 408
\parencite{en408}. Dort sind die Verfahren zur Bestimmung von
Biegefestigkeit und Biege-Elastizitaetsmodul im Vier-Punkt-Biegeversuch,
von Zug- und Druckfestigkeit parallel zur Faser sowie von Scherfestigkeit,
\gls{rohdichte} und Holzfeuchte festgelegt. Damit ist auch der formale Weg
beschrieben, auf dem Altholz ueberhaupt eine Festigkeitsklasse erhalten kann:
durch Neubewertung nach EN 384 auf der Grundlage von Pruefungen nach EN 408.
Dieser Weg ist es, den die vorliegende Arbeit methodisch nachbildet, wenn sie
\acrshort{mor} und \acrshort{moe} an Proben aus \gls{druckzone} und
\gls{zugzone} bestimmt.

\section{Bemessung nach Eurocode 5}
\label{sec:ec5}

Die Bemessung von Holztragwerken richtet sich nach Eurocode 5
\parencite{ec5}. Die Norm behandelt Vollholz, Brettschichtholz,
Holzwerkstoffe und Verbundkonstruktionen und deckt Tragfaehigkeit,
Gebrauchstauglichkeit und Dauerhaftigkeit ab. Fuer die Bemessung werden
charakteristische Werte fuer Biege-, Zug-, Druck- und Schubfestigkeit sowie
Elastizitaetsmoduln benoetigt, die entweder der Klassifizierung nach EN 338
entnommen oder nach EN 384 und EN 408 durch Pruefung ermittelt werden. Aus den
charakteristischen Werten entstehen Bemessungswerte, indem sie mit dem
Modifikationsbeiwert $k_{\mathrm{mod}}$ multipliziert und durch einen
Materialbeiwert geteilt werden. Der Modifikationsbeiwert beruecksichtigt die
Dauer der Lasteinwirkung und die klimatischen Bedingungen der
Nutzungsklasse.

Damit ist die Frage beantwortet, wie ein Bauteil ohne bekannte Herkunft
bemessen wird, allerdings nur formal: es muss zuvor in eine Festigkeitsklasse
eingeordnet werden, und diese Einordnung geschieht auf Grundlage dessen, was
sich am Bauteil heute beobachten oder messen laesst. Die Recherche fand in
Eurocode 5 keine eigenstaendigen Regelungen fuer Altholz. Praktisch bedeutet
das, dass ein wiederverwendeter Balken wie neues Holz bemessen wird, sobald er
eine Klasse erhalten hat, gegebenenfalls erweitert um einen Nachweis der
Dauerhaftigkeit gegen Pilz- und Insektenbefall und um eine Dokumentation
ausgeschlossener Kontaminationen. Ob eine solche Gleichbehandlung zutrifft,
ist eine empirische Frage, die die Norm nicht stellt. In der Praxis wird
reklassifiziertes Altholz nach Angaben der Fachliteratur eher in sekundaeren
tragenden Bauteilen wie Deckenbalken oder Dachsparren von
Sanierungsvorhaben eingesetzt als in primaeren Tragelementen.

\section{Die normative Luecke}
\label{sec:luecke}

Fasst man die Regelwerke zusammen, ergibt sich ein geschlossener Weg vom
Abfall zum Bauteil: die \acrshort{altholzv} entscheidet ueber die stoffliche
Zulaessigkeit, das \acrshort{krwg} raeumt der Wiederverwendung Vorrang ein,
DIN 4074 und EN 14081 sortieren, EN 338 und EN 1912 klassifizieren, EN 384 und
EN 408 liefern die charakteristischen Werte, und Eurocode 5 bemisst. An keiner
Stelle dieses Weges wird nach der Lastgeschichte des Bauteils gefragt.

Diese Luecke ist keine Nachlaessigkeit, sondern eine Folge der Bauart des
Systems. Sortierung und Klassifizierung sind fuer die Erstverwendung von
frischem Schnittholz entwickelt worden. Sie charakterisieren eine Population
und ordnen dem einzelnen Stueck eine Klasse anhand von Merkmalen zu, die im
Moment der Sortierung sichtbar oder messbar sind. Eine Vorbelastung ist
naturgemaess keines dieser Merkmale: sie hinterlaesst keine Astigkeit und
keine Faserneigung, sondern vermutlich nur Veraenderungen im Gefuege, die
weder die visuelle Beurteilung nach DIN 4074 noch die maschinelle Messung nach
EN 14081 erfassen kann. Der Modifikationsbeiwert $k_{\mathrm{mod}}$ zeigt, dass
die Norm die Bedeutung der Lastdauer durchaus kennt, richtet diese
Beruecksichtigung aber ausschliesslich in die Zukunft: er bewertet die
Einwirkungsdauer der kuenftigen Belastung im Bemessungszeitraum. Ein
Gegenstueck fuer bereits getragene Last existiert in der betrachteten
Normenkette nicht.

Fuer ein wiederverwendetes Bauteil hat das eine unmittelbare Konsequenz. Ein
Balken, der jahrzehntelang auf Biegung belastet war, ist kein homogener
Koerper: seine frueheren Zug- und Druckbereiche haben unterschiedliche
Spannungszustaende erfahren und liegen im reklassifizierten Bauteil dennoch
unter derselben Festigkeitsklasse. Ob diese Gleichbehandlung gerechtfertigt
ist, ist bislang nicht geklaert. Die vorliegende Arbeit setzt an dieser Stelle
an und prueft experimentell, ob sich die Restfestigkeit zwischen der
frueheren Druckzone und der frueheren Zugzone eines gebrauchten Balkens
messbar unterscheidet. Kapitel~\ref{ch:literature} sichtet dazu zunaechst den
Forschungsstand zu Altholzfestigkeit und Lastdauereffekten.
